\def\handoutsay{false}\def\handoutfill{false}
\documentclass[12pt]{article}
\newcommand{\handoutrunning}{Lecture 5: Block Matrices and Determinants}
\providecommand{\handoutsay}{true}
\usepackage[T1]{fontenc}
\usepackage{lmodern,microtype}
\usepackage[a4paper,margin=15mm,headheight=14pt,headsep=5mm]{geometry}
\usepackage{amsmath,amssymb,amsthm,mathtools}
\usepackage{enumitem,booktabs,needspace}
\usepackage{pifont}
\usepackage[dvipsnames]{xcolor}
\usepackage{fancyhdr}
\usepackage{ifthen}
\usepackage[hidelinks]{hyperref}
\hypersetup{pdfauthor={Chang-Ye Tu}}
\pagestyle{fancy}
\fancyhf{}
\fancyhead[L]{\small Linear Algebra}
\fancyhead[R]{\small\handoutrunning}
\fancyfoot[C]{\small\thepage}
\renewcommand{\headrulewidth}{0.3pt}
\setlength{\parindent}{0pt}
\setlength{\parskip}{4pt plus 1pt minus 1pt}
\setlength{\emergencystretch}{1.5em}
\usepackage{etoolbox}
\newcommand{\handoutdisplayskips}{\setlength{\abovedisplayskip}{7pt plus 2pt minus 2pt}%
  \setlength{\belowdisplayskip}{7pt plus 2pt minus 2pt}}
\apptocmd{\normalsize}{\handoutdisplayskips}{}{\PackageWarning{handout}{could not patch normalsize}}
\handoutdisplayskips
\setlist[enumerate]{label=(\alph*),leftmargin=*,itemsep=2pt,topsep=3pt}
\setlist[itemize]{leftmargin=*,itemsep=2pt,topsep=3pt}
\newcommand{\R}{\mathbb R}
\newcommand{\C}{\mathbb C}
\newcommand{\iu}{\mathrm{i}}
\newcommand{\contradiction}{\mathord{\text{\large\ding{56}}}}
\newcommand{\one}{\mathbf 1}
\DeclareMathOperator{\tr}{tr}
\DeclareMathOperator{\col}{col}
\DeclareMathOperator{\diag}{diag}
\DeclareMathOperator{\rank}{rank}
\DeclareMathOperator\prb{\mathsf{P}}
\DeclareMathOperator\expc{\mathsf{E}}
\DeclareMathOperator\var{var}
\DeclareMathOperator\cov{cov}
\DeclareMathOperator\cor{cor}
\DeclareMathOperator\argmin{argmin}
\DeclareMathOperator\argmax{argmax}
\DeclareMathOperator{\Span}{span}
\theoremstyle{definition}
\newtheorem{theorem}{Theorem}
\newtheorem{proposition}[theorem]{Proposition}
\newtheorem{lemma}[theorem]{Lemma}
\newtheorem{corollary}[theorem]{Corollary}
\newtheorem{definition}[theorem]{Definition}
\newtheorem{example}[theorem]{Example}
\newtheorem{remark}[theorem]{Remark}
\newtheorem{fact}{Fact}
\renewcommand{\thefact}{F\arabic{fact}}
\AddToHook{cmd/subsection/before}{\par\Needspace{9\baselineskip}}
\AddToHook{env/theorem/before}{\par\Needspace{4\baselineskip}}
\AddToHook{env/proposition/before}{\par\Needspace{4\baselineskip}}
\AddToHook{env/lemma/before}{\par\Needspace{4\baselineskip}}
\AddToHook{env/definition/before}{\par\Needspace{4\baselineskip}}
\AddToHook{env/example/before}{\par\Needspace{4\baselineskip}}
\definecolor{saybg}{gray}{0.86}
\NewDocumentEnvironment{say}{+b}
  {\ifthenelse{\equal{\handoutsay}{true}}{\par\Needspace{2\baselineskip}\noindent\colorbox{saybg}{\parbox{\dimexpr\linewidth-2\fboxsep\relax}%
   {\small\sffamily\ding{43}\enspace\textbf{Say.}\enspace#1}}\par}{}}
  {}
\providecommand{\handoutfill}{false}
\newsavebox{\boardbox}
\newlength{\boardht}
\definecolor{boardrule}{gray}{0.55}
\makeatletter
\newcommand{\boardstart}{\par\if@inlabel\leavevmode\par\fi}
\makeatother
\NewDocumentEnvironment{board}{o +b}
  {\boardstart\IfValueT{#1}{\noindent#1\par\nopagebreak}%
   \noindent\begin{lrbox}{\boardbox}\begin{minipage}{\linewidth}#2\end{minipage}\end{lrbox}%
   \setlength{\boardht}{\dimexpr\ht\boardbox+\dp\boardbox\relax}%
   \ifthenelse{\equal{\handoutfill}{true}}%
     {\setlength{\boardht}{1.4\boardht}%
      {\setlength{\fboxsep}{0pt}\setlength{\fboxrule}{0.4pt}\color{boardrule}%
       \framebox[\dimexpr\linewidth-2\fboxrule\relax]{\rule{0pt}{\boardht}}}}%
     {\ifthenelse{\equal{\handoutsay}{true}}%
        {\llap{\color{boardrule}\rule[-\dp\boardbox]{1.5pt}{\boardht}\hspace{4pt}}}{}%
      \usebox{\boardbox}}%
   \par}
  {}
\newcommand{\blank}[1]{\ifthenelse{\equal{\handoutfill}{true}}{\underline{\phantom{#1}}}{#1}}
\newcommand{\lec}[1]{Lecture~#1}
\newcommand{\handouttitle}[2]{%
  \begin{center}
  {\LARGE\sffamily\bfseries #2\par}\vspace{3pt}
  {\sffamily Linear Algebra \textperiodcentered\ \lec{#1}\par}
  \ifthenelse{\equal{\handoutsay}{true}}{\vspace{3pt}{\small\sffamily\colorbox{saybg}{\ding{43}\enspace\textbf{Say.}}\enspace
    teacher's copy: the shaded boxes are not in the students' copy.\par}}{}
  \ifthenelse{\equal{\handoutfill}{true}}{\vspace{3pt}{\small\sffamily Fill-in copy: complete the framed spaces in
    class.\par}}{}
  \end{center}}

\hypersetup{pdftitle={Linear Algebra, Lecture 5: Block Matrices and Determinants}}
\newcommand{\boardtop}{\vspace*{-\abovedisplayskip}}
\DeclareMathOperator{\row}{row}
\DeclareMathOperator{\adj}{adj}
\begin{document}
\handouttitle{5}{Block Matrices and Determinants}

\textbf{Last time} (\lec{4}, Definition~10, Theorem~12(d)): \(\rank A\) is the number of
nonzero rows of an RREF of \(A\); \(A\in\R^{n\times n}\) is invertible \(\iff\) \(\rank A=n\). Throughout, \(e_j\) is
column \(j\) of \(I\); \(\row_i\), \(\col_j\): row \(i\), column \(j\); \(A(i|j)\): \(A\) without row \(i\) and column \(j\).

\subsection*{1. Block matrices}

\begin{definition}\label{b5.1}\label{b5.2}
A \textbf{partition} of \(A\in\R^{m\times n}\) splits the rows into consecutive groups of \(m_1,\dots,m_r\) rows and
the columns into consecutive groups of \(n_1,\dots,n_s\) columns. With \(M_0=N_0=0\), \(M_i=m_1+\dots+m_i\),
\(N_j=n_1+\dots+n_j\), the \textbf{block} \(A_{ij}\in\R^{m_i\times n_j}\) has the entries
\((A_{ij})_{xy}=a_{M_{i-1}+x,\,N_{j-1}+y}\); \(A=(A_{ij})\) is the array of its blocks. Matrices of the same size,
partitioned alike, are equal \(\iff\) their blocks are equal. A partition of \(B\in\R^{m\times n}\) is \textbf{conformal
for addition} if it has the same groups; one of \(B\in\R^{n\times p}\) is \textbf{conformal for multiplication} (in
\(AB\)) if its row groups are \(n_1,\dots,n_s\) (its column groups \(p_1,\dots,p_t\) are free).
\end{definition}

\Needspace{15\baselineskip}
\begin{theorem}\label{b5.3}\label{b5.6}
Let \(A\in\R^{m\times n}\) be partitioned into blocks \(A_{ij}\in\R^{m_i\times n_j}\) (\(1\leqslant i\leqslant r\),
\(1\leqslant j\leqslant s\)).
\begin{enumerate}
\item If \(B\in\R^{m\times n}\) is partitioned conformally for addition, into blocks \(B_{ij}\), then \(A+B\) has the blocks
\(A_{ij}+B_{ij}\), and \(cA\) the blocks \(cA_{ij}\) (\(c\in\R\)).
\item \(A'\), with row groups \(n_1,\dots,n_s\) and column groups \(m_1,\dots,m_r\), has the \((j,i)\) block \(A_{ij}'\).
\item If \(B\in\R^{n\times p}\) is partitioned conformally for multiplication, into blocks
\(B_{jk}\in\R^{n_j\times p_k}\), then \(AB\), with row groups \(m_1,\dots,m_r\) and column groups \(p_1,\dots,p_t\), has the
blocks \((AB)_{ik}=A_{i1}B_{1k}+\dots+A_{is}B_{sk}\in\R^{m_i\times p_k}\).
\end{enumerate}
\end{theorem}
\begin{proof}
(a) For \(u=M_{i-1}+x\), \(v=N_{j-1}+y\): \((A+B)_{uv}=a_{uv}+b_{uv}=(A_{ij}+B_{ij})_{xy}\), \((cA)_{uv}=(cA_{ij})_{xy}\)
(\lec{1}, Definition~1).
\begin{board}[(b) Row group \(j\) of \(A'\): rows \(N_{j-1}+1,\dots,N_j\); column group \(i\): columns
\(M_{i-1}+1,\dots,M_i\) (\lec{1}, Definition~7):]
\[
(A')_{N_{j-1}+y,\,M_{i-1}+x}=a_{M_{i-1}+x,\,N_{j-1}+y}=(A_{ij})_{xy}=(A_{ij}')_{yx}
\]
\end{board}
\Needspace{3\baselineskip}
(c)
\begin{itemize}
\item \(A_{i*}\in\R^{m_i\times n}\): the rows of \(A\) in row group \(i\), \(A_{i*}=[A_{i1}\ \cdots\ A_{is}]\);
\(B_{*k}\in\R^{n\times p_k}\): the columns of \(B\) in column group \(k\), the blocks \(B_{1k},\dots,B_{sk}\) stacked
(conformality).
\item The columns of \(AB\) in group \(k\) are those of \(AB_{*k}\) (\(\col_j(AB)=Ab_j\)); the rows of \(AB_{*k}\) in group
\(i\) are those of \(A_{i*}B_{*k}\) (\(\row_i(AB)=\row_i(A)B\)) (\lec{1}, Theorem~4(a),(c)). So
\((AB)_{ik}=A_{i*}B_{*k}\).
\item Inner sum: \(A_{i*}B_{*k}=\sum_{l=1}^n\col_l(A_{i*})\row_l(B_{*k})\) by the outer product expansion
(\lec{1}, Theorem~4(d)), and the indices \(N_{j-1}<l\leqslant N_j\) give
\(\sum_{y=1}^{n_j}\col_y(A_{ij})\row_y(B_{jk})=A_{ij}B_{jk}\); adding the \(s\) stretches,
\((AB)_{ik}=\sum_jA_{ij}B_{jk}\).\qedhere
\end{itemize}
\end{proof}
For \(r=s=t=2\):
\[
\begin{pmatrix}A_{11}&A_{12}\\A_{21}&A_{22}\end{pmatrix}\begin{pmatrix}B_{11}&B_{12}\\B_{21}&B_{22}\end{pmatrix}
=\begin{pmatrix}A_{11}B_{11}+A_{12}B_{21}&A_{11}B_{12}+A_{12}B_{22}\\A_{21}B_{11}+A_{22}B_{21}&A_{21}B_{12}+A_{22}B_{22}\end{pmatrix}
\]
\(\diag(A_1,\dots,A_r)\) has the square diagonal blocks \(A_i\) and zero blocks elsewhere; by (c), with the same groups,
\(\diag(A_1,\dots,A_r)\diag(B_1,\dots,B_r)=\diag(A_1B_1,\dots,A_rB_r)\) (only \(j=i=k\) gives a nonzero term).

\subsection*{2. Block elimination and the Schur complement}

\Needspace{12\baselineskip}
\begin{theorem}\label{b5.9}\label{b5.8}
Let \(M=\begin{pmatrix}A&B\\C&D\end{pmatrix}\) with \(A\in\R^{p\times p}\) invertible, \(B\in\R^{p\times q}\),
\(C\in\R^{n\times p}\), \(D\in\R^{n\times q}\), and \(S=D-CA^{-1}B\in\R^{n\times q}\), the \textbf{Schur complement} of
\(A\) in \(M\). Then
\[
\begin{pmatrix}I_p&0\\-CA^{-1}&I_n\end{pmatrix}M=\begin{pmatrix}A&B\\0&S\end{pmatrix},\qquad
M\begin{pmatrix}I_p&-A^{-1}B\\0&I_q\end{pmatrix}=\begin{pmatrix}A&0\\C&S\end{pmatrix}
\]
\end{theorem}
\begin{proof}
\begin{board}[Left factor: column groups \((p,n)\), the row groups of \(M\) (Theorem~\ref{b5.3}(c); \lec{1},
Theorem~3(a)):]
\[
\begin{pmatrix}I_p&0\\-CA^{-1}&I_n\end{pmatrix}M
=\begin{pmatrix}I_pA+0C&I_pB+0D\\-CA^{-1}A+I_nC&-CA^{-1}B+I_nD\end{pmatrix}=\begin{pmatrix}A&B\\0&D-CA^{-1}B\end{pmatrix}
\]
\end{board}
\begin{board}[Right factor: row groups \((p,q)\), the column groups of \(M\) (Theorem~\ref{b5.3}(c)):]
\[
M\begin{pmatrix}I_p&-A^{-1}B\\0&I_q\end{pmatrix}
=\begin{pmatrix}AI_p+B0&-AA^{-1}B+BI_q\\CI_p+D0&-CA^{-1}B+DI_q\end{pmatrix}=\begin{pmatrix}A&0\\C&D-CA^{-1}B\end{pmatrix}
\]
\end{board}
\nopagebreak
\end{proof}
When \(D\) is invertible, \(A-BD^{-1}C\) is likewise the \textbf{Schur complement} of \(D\) in \(M\).

\Needspace{14\baselineskip}
\begin{corollary}\label{b5.10}\label{b5.12}
Let \(M=\begin{pmatrix}A&B\\C&D\end{pmatrix}\) with \(A\in\R^{p\times p}\), \(D\in\R^{q\times q}\), and let \(A\) and
\(S=D-CA^{-1}B\) be invertible. Then \(M\) is invertible and
\[
M^{-1}=\begin{pmatrix}I_p&-A^{-1}B\\0&I_q\end{pmatrix}\begin{pmatrix}A^{-1}&0\\0&S^{-1}\end{pmatrix}
\begin{pmatrix}I_p&0\\-CA^{-1}&I_q\end{pmatrix}
=\begin{pmatrix}A^{-1}+A^{-1}BS^{-1}CA^{-1}&-A^{-1}BS^{-1}\\-S^{-1}CA^{-1}&S^{-1}\end{pmatrix}
\]
\end{corollary}
\begin{proof}
\begin{board}[\(E_3\), \(E_2\), \(E_1\): the three factors, from left to right; Theorem~\ref{b5.3}(c):]
\boardtop
\begin{align*}
E_1M&=\begin{pmatrix}A&B\\0&S\end{pmatrix},\quad
E_2E_1M=\begin{pmatrix}A^{-1}A&A^{-1}B\\S^{-1}\cdot0&S^{-1}S\end{pmatrix}=\begin{pmatrix}I_p&A^{-1}B\\0&I_q\end{pmatrix}
&&\text{(Theorem~\ref{b5.9}, \(n=q\))}\\
E_3E_2E_1M&=\begin{pmatrix}I_p&A^{-1}B-A^{-1}B\\0&I_q\end{pmatrix}=I\implies M^{-1}=E_3E_2E_1
&&\text{(\lec{2}, Theorem~15)}
\end{align*}
\end{board}
\begin{board}[Multiplying out (Theorem~\ref{b5.3}(c)):]
\[
E_2E_1=\begin{pmatrix}A^{-1}&0\\-S^{-1}CA^{-1}&S^{-1}\end{pmatrix},\qquad
E_3(E_2E_1)=\begin{pmatrix}A^{-1}+A^{-1}BS^{-1}CA^{-1}&-A^{-1}BS^{-1}\\-S^{-1}CA^{-1}&S^{-1}\end{pmatrix}
\]
\end{board}
\nopagebreak
\end{proof}
On \([M\mid I]\) with block rows \(R_1,R_2\), left multiplication by \(E_1\), \(E_2\), \(E_3\) is the \textbf{block row
operation} \(R_2\leftarrow R_2-CA^{-1}R_1\), then \(R_1\leftarrow A^{-1}R_1\) and \(R_2\leftarrow S^{-1}R_2\), then
\(R_1\leftarrow R_1-A^{-1}BR_2\) (Theorem~\ref{b5.3}(c)); together they give \([I\mid M^{-1}]\): \textbf{block elimination}.

\Needspace{16\baselineskip}
\begin{example}\label{b5.11}
Find the Schur complement \(S\) of \(A\), and \(M^{-1}\), for \(M\in\R^{5\times5}\) with groups \((2,3)\) for rows and
columns:
\[
M=\begin{pmatrix}A&B\\C&D\end{pmatrix},\quad A=\begin{pmatrix}2&1\\1&1\end{pmatrix},\quad
B=\begin{pmatrix}2&1&3\\1&1&2\end{pmatrix},\quad C=B',\quad D=\begin{pmatrix}3&2&3\\2&3&3\\3&3&7\end{pmatrix}
\]
\end{example}
\textit{Solution.}
\begin{board}[\(A^{-1}\) (\lec{2}, Proposition~17: \(2\cdot1-1\cdot1=1\)); \(A^{-1}B\), \(CA^{-1}\); \(S\) (Theorem~\ref{b5.9}):]
\[
\begin{gathered}
A^{-1}=\begin{pmatrix}1&-1\\-1&2\end{pmatrix},\qquad
A^{-1}B=\begin{pmatrix}1&0&1\\0&1&1\end{pmatrix},\qquad
CA^{-1}=\begin{pmatrix}2&1\\1&1\\3&2\end{pmatrix}\begin{pmatrix}1&-1\\-1&2\end{pmatrix}=\begin{pmatrix}1&0\\0&1\\1&1\end{pmatrix}\\
S=D-CA^{-1}B=\begin{pmatrix}3&2&3\\2&3&3\\3&3&7\end{pmatrix}-\begin{pmatrix}2&1&3\\1&1&2\\3&2&5\end{pmatrix}
=\begin{pmatrix}1&1&0\\1&2&1\\0&1&2\end{pmatrix}
\end{gathered}
\]
\end{board}
\begin{board}[\(S^{-1}\) (\(SS^{-1}=I_3\); \lec{2}, Theorem~15); the off-diagonal blocks of \(M^{-1}\)
(Corollary~\ref{b5.10}):]
\[
S^{-1}=\begin{pmatrix}3&-2&1\\-2&2&-1\\1&-1&1\end{pmatrix},\qquad
-A^{-1}BS^{-1}=\begin{pmatrix}-4&3&-2\\1&-1&0\end{pmatrix},\qquad
-S^{-1}CA^{-1}=\begin{pmatrix}-4&1\\3&-1\\-2&0\end{pmatrix}
\]
\end{board}
\begin{board}[The upper left block \(A^{-1}+A^{-1}BS^{-1}CA^{-1}=A^{-1}-(A^{-1}B)(-S^{-1}CA^{-1})\); \(M^{-1}\)
(Corollary~\ref{b5.10}); a check of row \(1\):]
\[
\begin{gathered}
\begin{pmatrix}1&-1\\-1&2\end{pmatrix}
-\begin{pmatrix}1&0&1\\0&1&1\end{pmatrix}\begin{pmatrix}-4&1\\3&-1\\-2&0\end{pmatrix}
=\begin{pmatrix}1&-1\\-1&2\end{pmatrix}-\begin{pmatrix}-6&1\\1&-1\end{pmatrix}=\begin{pmatrix}7&-2\\-2&3\end{pmatrix}\\
M^{-1}=\left(\begin{array}{cc|ccc}7&-2&-4&3&-2\\-2&3&1&-1&0\\\hline-4&1&3&-2&1\\3&-1&-2&2&-1\\-2&0&1&-1&1\end{array}\right),\qquad
\begin{aligned}\row_1(M)M^{-1}&=2(7,-2,-4,3,-2)+(-2,3,1,-1,0)\\&\quad+2(-4,1,3,-2,1)+(3,-1,-2,2,-1)\\
&\quad+3(-2,0,1,-1,1)\\&=(1,0,0,0,0)\end{aligned}
\end{gathered}
\]
\end{board}

\Needspace{8\baselineskip}
\begin{theorem}\label{b5.12b}
Let \(M=\begin{pmatrix}A&B\\C&D\end{pmatrix}\) with \(A\in\R^{p\times p}\), \(D\in\R^{q\times q}\), and let \(D\) and
\(F=A-BD^{-1}C\) be invertible. Then \(M\) is invertible and
\[
M^{-1}=\begin{pmatrix}F^{-1}&-F^{-1}BD^{-1}\\-D^{-1}CF^{-1}&D^{-1}+D^{-1}CF^{-1}BD^{-1}\end{pmatrix}
\]
\end{theorem}
\begin{proof}
For \(X=\left(\begin{smallmatrix}P&Q\\R&T\end{smallmatrix}\right)\), by Theorem~\ref{b5.3}(c) and
Definition~\ref{b5.1}, \(MX=I\) \(\iff\) \(AP+BR=I_p\), \(CP+DR=0\), \(AQ+BT=0\), \(CQ+DT=I_q\).
\begin{board}[Solve the four equations (\lec{1}, Theorem~3(a)):]
\boardtop
\begin{align*}
&R=-D^{-1}CP\implies(A-BD^{-1}C)P=FP=I_p\implies P=F^{-1},\ R=-D^{-1}CF^{-1}\\
&T=D^{-1}-D^{-1}CQ\implies AQ+BT=FQ+BD^{-1}=0\\
&\implies Q=-F^{-1}BD^{-1},\ T=D^{-1}+D^{-1}CF^{-1}BD^{-1}
\end{align*}
\end{board}
\begin{board}[Conversely, these blocks satisfy the four equations:]
\boardtop
\begin{align*}
AP+BR&=(A-BD^{-1}C)F^{-1}=I_p,\qquad CP+DR=CF^{-1}-CF^{-1}=0\\
AQ+BT&=-(A-BD^{-1}C)F^{-1}BD^{-1}+BD^{-1}=0\\
CQ+DT&=-CF^{-1}BD^{-1}+I_q+CF^{-1}BD^{-1}=I_q\\
MX&=I\implies M^{-1}=X\qquad\text{(\lec{2}, Theorem~15)}
\end{align*}
\end{board}
\nopagebreak
\end{proof}

\subsection*{3. Determinants}

\begin{definition}\label{b5.20}
For \(A\in\R^{n\times n}\), the \textbf{determinant} \(\det A\) is defined by induction on \(n\): \(\det(a)=a\) for
\(n=1\); for \(n\geqslant2\), with the \textbf{minors} \(\det A(i|j)\) and the \textbf{cofactors}
\(c_{ij}=(-1)^{i+j}\det A(i|j)\),
\[
\det A=a_{11}c_{11}+a_{21}c_{21}+\dots+a_{n1}c_{n1}\qquad\text{(expansion along column \(1\))}
\]
For \(1\leqslant j\leqslant n\) let \(E_j(A)=a_{1j}c_{1j}+\dots+a_{nj}c_{nj}\) (along column \(j\)), and \(E_1(a)=a\) for
\(n=1\). So \(E_1(A)=\det A\).
\end{definition}

\Needspace{7\baselineskip}
\begin{theorem}\label{b5.13}
Let \(n\geqslant1\), \(1\leqslant j\leqslant n\), \(A\in\R^{n\times n}\), rows \(r\ne s\), and \(c\in\R\).
\begin{enumerate}
\item \(E_j\) is linear in each row: multiplying row \(r\) by \(c\) multiplies \(E_j(A)\) by \(c\); if \(A_1,A_2,A_3\) agree
except in row \(r\) and \(\row_r(A_3)=\row_r(A_1)+\row_r(A_2)\), then \(E_j(A_3)=E_j(A_1)+E_j(A_2)\).
\item Adding \(c\) times row \(s\) to row \(r\) does not change \(E_j(A)\).
\item Interchanging rows \(r\) and \(s\) changes the sign of \(E_j(A)\). So \(E_j(A)=0\) if two rows of \(A\) are equal.
\item \(E_j(I_n)=1\).
\end{enumerate}
\end{theorem}
\begin{proof}
Induction on \(n\). For \(n=1\), \(E_1(a)=a\) satisfies (a) and (d); (b), (c) concern two rows. Let \(n\geqslant2\), assume
(a)--(d) for \(\det=E_1\) on \(\R^{(n-1)\times(n-1)}\), and write \(E=E_j\),
\(E(A)=\sum_i(-1)^{i+j}a_{ij}\det A(i|j)\).
\begin{itemize}
\item (a): \(i\ne r\): \(a_{ij}\) does not involve row \(r\), and \(\det A(i|j)\) is linear in it (\(A(i|j)\) contains row
\(r\) without entry \(j\); induction); \(i=r\): \(\det A(r|j)\) does not involve row \(r\), and \(a_{rj}\) is linear in it.
\item Rows \(k\), \(k+1\) equal \(\Rightarrow\) \(E(A)=0\): \(\det A(i|j)=0\) for \(i\ne k,k+1\) (two equal rows; induction,
(c)); the terms \(i=k,k+1\) cancel, as \(a_{kj}=a_{k+1,j}\), \(A(k|j)=A(k+1|j)\) and the signs are \((-1)^{k+j}\),
\((-1)^{k+1+j}\).
\item Adjacent interchange: with \(E(u,v)\) the value when rows \(k\), \(k+1\) are \(u'\), \(v'\) (\(u,v\in\R^n\)),
\(0=E(u+v,u+v)=E(u,u)+E(u,v)+E(v,u)+E(v,v)=E(u,v)+E(v,u)\) ((a), previous point).
\item (c): interchanging rows \(r<s\) is \(2(s-r)-1\) adjacent interchanges (row \(s\) up to position \(r\), then the former
row \(r\) down to position \(s\)), an odd number. Two equal rows: interchanging them, \(E(A)=-E(A)\).
\item (b): by (a) in row \(r\), the new value is \(E(A)+cE(A_s)\), \(A_s\) = \(A\) with row \(r\) replaced by row \(s\);
\(E(A_s)=0\) by (c).
\item (d): \(\col_j(I_n)=e_j\), \(I_n(j|j)=I_{n-1}\): \(E(I_n)=(-1)^{2j}\det I_{n-1}=1\) (induction).\qedhere
\end{itemize}
\end{proof}

\Needspace{13\baselineskip}
\begin{theorem}\label{b5.14}
Let \(A\in\R^{n\times n}\) and \(1\leqslant j\leqslant n\).
\begin{enumerate}
\item If \(A\) has a zero row, then \(E_j(A)=0\).
\item If \(A\) is upper triangular (\lec{1}, Definition~12), then \(E_j(A)=a_{11}a_{22}\cdots a_{nn}\); in particular
\(E_j(\diag(d_1,\dots,d_n))=d_1\cdots d_n\).
\item \(E_j(D_r(c))=c\), \(E_j(T_{rs}(c))=1\), \(E_j(P_{rs})=-1\), and \(E_j(FA)=E_j(F)E_j(A)\) for every elementary matrix
\(F\in\R^{n\times n}\) (\lec{2}, Definition~13).
\end{enumerate}
\end{theorem}
\begin{proof}
\begin{board}[(a) Multiplying the zero row by \(2\) does not change \(A\) (Theorem~\ref{b5.13}(a)):]
\[
E_j(A)=2E_j(A)\implies E_j(A)=0
\]
\end{board}
\Needspace{4\baselineskip}
(b)
\begin{itemize}
\item All \(a_{ii}\ne0\): \(R_i\leftarrow a_{ii}^{-1}R_i\) (all \(i\); \lec{2}, Definition~2) gives \(B\), upper triangular with diagonal entries
\(1\), \(E_j(B)=E_j(A)/(a_{11}\cdots a_{nn})\) (Theorem~\ref{b5.13}(a)); backward elimination (\lec{2}, Definition~5) with
these pivots turns \(B\) into \(I_n\) by additions: \(E_j(B)=E_j(I_n)=1\) (Theorem~\ref{b5.13}(b),(d)).
\item Some \(a_{ii}=0\), \(i\) the largest: row \(i\) is \(0\) in columns \(1,\dots,i\). For \(l=i+1,\dots,n\) in turn,
subtract a multiple of row \(l\) from row \(i\) to clear entry \((i,l)\) (\(a_{ll}\ne0\); row \(l\) is \(0\) before column
\(l\), so the zeros made stay). This keeps \(E_j(A)\) (Theorem~\ref{b5.13}(b)) and makes row \(i\) zero:
\(E_j(A)=0=a_{11}\cdots a_{nn}\) by (a).
\end{itemize}
\begin{board}[(c) \(D_r(c)\), \(T_{rs}(c)\), \(P_{rs}\) are \(I_n\) after \(R_r\leftarrow cR_r\),
\(R_r\leftarrow R_r+cR_s\), \(R_r\leftrightarrow R_s\) (\lec{2}, Definition~13):]
\(E_j(D_r(c))=cE_j(I_n)=c\), \(E_j(T_{rs}(c))=E_j(I_n)=1\), \(E_j(P_{rs})=-E_j(I_n)=-1\)
(Theorem~\ref{b5.13}).\\
\(FA\) is \(A\) after the operation of \(F\) (\lec{2}, Theorem~14(a)), so \(E_j(FA)\) is \(cE_j(A)\),
\(E_j(A)\) or \(-E_j(A)\) (Theorem~\ref{b5.13}): \(E_j(FA)=E_j(F)E_j(A)\).
\end{board}
\nopagebreak
\end{proof}

\Needspace{14\baselineskip}
\begin{theorem}\label{b5.15}
Let interchanges of two rows and additions of a multiple of one row to another take \(A\in\R^{n\times n}\) to an upper
triangular \(U\), with \(t\) interchanges. Then \(E_j(A)=(-1)^tu_{11}\cdots u_{nn}\) for every \(j\). Forward elimination
(Steps 1, 2, 4 of \lec{2}, Definition~5, without scaling) always ends with such a \(U\). Hence \(E_j(A)=\det A\) for every \(j\):
\[
\det A=a_{1j}c_{1j}+a_{2j}c_{2j}+\dots+a_{nj}c_{nj}\qquad(\textbf{expansion along column }j)
\]
\end{theorem}
\begin{proof}
\begin{board}[Each addition keeps \(E_j\), each interchange changes its sign (Theorem~\ref{b5.13}(b),(c)):]
\(u_{11}\cdots u_{nn}=E_j(U)=(-1)^tE_j(A)\) (Theorem~\ref{b5.14}(b)).\\
Forward elimination: leading entries in columns \(k_1<k_2<\cdots\), so \(k_i\geqslant i\), and zero rows at the bottom
(\lec{2}, Definition~5): \(U\) is upper triangular.\\
So \(E_j(A)=(-1)^tu_{11}\cdots u_{nn}=E_1(A)=\det A\) for every \(j\).
\end{board}
\nopagebreak
\end{proof}
If a row is multiplied by \(c\ne0\) on the way, \(\det\) is multiplied by \(c\) (Theorem~\ref{b5.13}(a)): divide by \(c\) at
the end.

\Needspace{19\baselineskip}
\begin{example}\label{b5.16}
Compute \(\det A\) for
\[
A=\begin{pmatrix}0&2&1&1\\1&2&1&0\\2&6&4&4\\1&0&2&3\end{pmatrix}
\]
\end{example}
\textit{Solution.}
\begin{board}[\(a_{11}=0\): begin with an interchange; then additions (Theorem~\ref{b5.15}):]
\[
\begin{gathered}
A\xrightarrow{R_1\leftrightarrow R_2}
\begin{pmatrix}1&2&1&0\\0&2&1&1\\2&6&4&4\\1&0&2&3\end{pmatrix}
\xrightarrow[R_4\leftarrow R_4-R_1]{R_3\leftarrow R_3-2R_1}
\begin{pmatrix}1&2&1&0\\0&2&1&1\\0&2&2&4\\0&-2&1&3\end{pmatrix}\\
\xrightarrow[R_4\leftarrow R_4+R_2]{R_3\leftarrow R_3-R_2}
\begin{pmatrix}1&2&1&0\\0&2&1&1\\0&0&1&3\\0&0&2&4\end{pmatrix}
\xrightarrow{R_4\leftarrow R_4-2R_3}
\begin{pmatrix}1&2&1&0\\0&2&1&1\\0&0&1&3\\0&0&0&-2\end{pmatrix}=U,\quad
t=1:\ \det A=-\bigl(1\cdot2\cdot1\cdot(-2)\bigr)=4
\end{gathered}
\]
\end{board}

\Needspace{16\baselineskip}
\begin{theorem}\label{b5.18}
Let \(A,B\in\R^{n\times n}\).
\begin{enumerate}
\item \(A\) is invertible \(\iff\) \(\det A\ne0\) \(\iff\) \(\rank A=n\).
\item \(\det(AB)=\det A\det B\).
\item If \(A\) is invertible, \(\det A^{-1}=1/\det A\). Moreover \(\det(cA)=c^n\det A\) (\(c\in\R\)) and
\(\det A^k=(\det A)^k\) (\(k=1,2,\dots\)).
\end{enumerate}
\end{theorem}
\begin{proof}
\begin{board}[(a) An RREF \(R=F_k\cdots F_1A\) of \(A\) (\lec{2}, Theorem~14(c)); \(\rank A\) is the number of nonzero
rows of \(R\) (\lec{4}, Definition~10); \(\det I_n=1\) (Theorem~\ref{b5.13}(d)):]
\boardtop
\begin{align*}
\det R&=\det F_k\cdots\det F_1\det A,\quad\det F_i\in\{c\ne0,\ 1,\ -1\}&&\text{(Theorem~\ref{b5.14}(c))}\\
A\text{ invertible}&\implies R=I_n\implies\det A\ne0&&\text{(\lec{2}, Theorem~15)}\\
A\text{ not invertible}&\iff\rank A<n\implies R\text{ has a zero row}&&\text{(\lec{4}, Theorem~12(d))}\\
&\implies\det R=0\implies\det A=0&&\text{(Theorem~\ref{b5.14}(a))}
\end{align*}
\end{board}
\begin{board}[(b) Two cases (\lec{2}, Theorem~15):]
\(A\) not invertible: then \(AB\) is not invertible (an inverse \(G\) of \(AB\) would give \(A(BG)=I\)), and
\(\det(AB)=0=\det A\det B\) by (a).\\
\(A\) invertible: \(A=F_1\cdots F_k\) with elementary \(F_i\), and \(\det(AB)=\det F_1\cdots\det F_k\det B\)
(Theorem~\ref{b5.14}(c)); with \(B=I_n\), \(\det A=\det F_1\cdots\det F_k\); so \(\det(AB)=\det A\det B\).
\end{board}
\begin{board}[(c) From (b); powers: \(A^{k+1}=A^kA\) (\lec{1}, after Theorem~3):]
\boardtop
\begin{align*}
\det A\det A^{-1}&=\det(AA^{-1})=\det I_n=1\\
\det(cA)&=\det(cI_n)\det A=c^n\det A&&\text{(Theorem~\ref{b5.14}(b))}\\
\det A^{k+1}&=\det(A^kA)=\det A^k\det A&&\text{(induction on \(k\))}
\end{align*}
\end{board}
\nopagebreak
\end{proof}

\Needspace{11\baselineskip}
\begin{theorem}\label{b5.21}\label{b5.22}
Let \(A\in\R^{n\times n}\).
\begin{enumerate}
\item If \(n\geqslant2\): \(\det A=a_{i1}c_{i1}+a_{i2}c_{i2}+\dots+a_{in}c_{in}\) for every \(i\) (\textbf{expansion along row}
\(i\)).
\item \(\det A'=\det A\).
\end{enumerate}
\end{theorem}
\begin{proof}
\begin{board}[(a) \(\row_i(A)=\sum_ja_{ij}e_j'\); \(G_j\) = \(A\) with row \(i\) replaced by \(e_j'\); \(H_j\) =
\(G_j\) after \(R_k\leftarrow R_k-a_{kj}R_i\) (\(k\ne i\)), so \(\col_j(H_j)=e_i\), \(H_j(i|j)=A(i|j)\):]
\boardtop
\begin{align*}
\det A&=\textstyle\sum_ja_{ij}\det G_j&&\text{(Theorem~\ref{b5.13}(a))}\\
\det G_j&=\det H_j=(-1)^{i+j}\det A(i|j)=c_{ij}&&\text{(Theorems~\ref{b5.13}(b), \ref{b5.15}, column \(j\))}
\end{align*}
\end{board}
\begin{board}[(b) Induction on \(n\); \(n=1\): \(A'=A\). For \(n\geqslant2\), \(A'(k|1)=A(1|k)'\):]
\boardtop
\begin{align*}
\det A'&=\textstyle\sum_k(-1)^{k+1}a_{1k}\det A(1|k)'&&\text{(Definition~\ref{b5.20})}\\
&=\textstyle\sum_k(-1)^{1+k}a_{1k}\det A(1|k)=\det A&&\text{(induction; (a), \(i=1\))}
\end{align*}
\end{board}
\nopagebreak
\end{proof}
An operation on the columns of \(A\) is the same operation on the rows of \(A'\); so Theorem~\ref{b5.13}(a)--(c) and
Theorem~\ref{b5.14}(a) hold for columns, and a lower triangular \(A\) has \(\det A=\det A'=a_{11}\cdots a_{nn}\)
(Theorem~\ref{b5.14}(b)). For \(n=2\), \(\det\left(\begin{smallmatrix}a&b\\c&d\end{smallmatrix}\right)=ad-bc\) (Definition~\ref{b5.20}), the number of
\lec{2}, Proposition~17.

\subsection*{4. The adjugate and Cramer's rule}

\Needspace{12\baselineskip}
\begin{theorem}\label{b5.24}\label{b5.25}
Let \(n\geqslant2\) and \(A\in\R^{n\times n}\) with cofactors \(c_{ij}\). The \textbf{adjugate} of \(A\) is
\(\adj A=(c_{ij})'\), that is, \((\adj A)_{ij}=c_{ji}\). Then \(A\adj A=(\det A)I_n\); if \(\det A\ne0\), then
\(A^{-1}=\frac1{\det A}\adj A\). If \(A\) is symmetric, then \(c_{ji}=c_{ij}\).
\end{theorem}
\begin{proof}
\begin{board}[Entry \((i,k)\) of \(A\adj A\); \(\tilde A\) = \(A\) with row \(k\) replaced by row \(i\) (\(k\ne i\)):]
\boardtop
\begin{align*}
(A\adj A)_{ik}&=\textstyle\sum_ja_{ij}c_{kj}=\det A\ (k=i)&&\text{(Theorem~\ref{b5.21}(a), row \(i\))}\\
\tilde A(k|j)=A(k|j)&\implies\textstyle\sum_ja_{ij}c_{kj}=\det\tilde A=0\ (k\ne i)
&&\text{(Theorem~\ref{b5.21}(a), row \(k\); Theorem~\ref{b5.13}(c))}\\
A\bigl(\tfrac1{\det A}\adj A\bigr)&=I_n\implies A^{-1}=\tfrac1{\det A}\adj A&&\text{(\lec{2}, Theorem~15)}
\end{align*}
\end{board}
\begin{board}[\(A\) symmetric (\(A=A'\); \lec{1}, Definition~7): delete row \(j\) and column \(i\):]
\[
A(j|i)=A'(j|i)=A(i|j)'\implies\det A(j|i)=\det A(i|j)\implies c_{ji}=c_{ij}\qquad\text{(Theorem~\ref{b5.21}(b))}
\]
\end{board}
\nopagebreak
\end{proof}

\Needspace{13\baselineskip}
\begin{example}\label{b5.26}
Find \(A^{-1}\) by Theorem~\ref{b5.24} for
\(A=\left(\begin{smallmatrix}1&1&0&0\\0&1&1&0\\0&0&1&1\\-1&0&0&1\end{smallmatrix}\right)\).
\end{example}
\textit{Solution.}
\begin{board}[Each \(3\times3\) minor is triangular (Theorem~\ref{b5.14}(b) and the lines after Theorem~\ref{b5.21}), or
has a row with one nonzero entry (Theorem~\ref{b5.21}(a)):]
\thickmuskip=2mu plus 1mu
\boardtop
\begin{alignat*}{4}
c_{11}&=\det\left(\begin{smallmatrix}1&1&0\\0&1&1\\0&0&1\end{smallmatrix}\right)=1,\ &
c_{12}&=-\det\left(\begin{smallmatrix}0&1&0\\0&1&1\\-1&0&1\end{smallmatrix}\right)=1,\ &
c_{13}&=\det\left(\begin{smallmatrix}0&1&0\\0&0&1\\-1&0&1\end{smallmatrix}\right)=-1,\ &
c_{14}&=-\det\left(\begin{smallmatrix}0&1&1\\0&0&1\\-1&0&0\end{smallmatrix}\right)=1\\
c_{21}&=-\det\left(\begin{smallmatrix}1&0&0\\0&1&1\\0&0&1\end{smallmatrix}\right)=-1,\ &
c_{22}&=\det\left(\begin{smallmatrix}1&0&0\\0&1&1\\-1&0&1\end{smallmatrix}\right)=1,\ &
c_{23}&=-\det\left(\begin{smallmatrix}1&1&0\\0&0&1\\-1&0&1\end{smallmatrix}\right)=1,\ &
c_{24}&=\det\left(\begin{smallmatrix}1&1&0\\0&0&1\\-1&0&0\end{smallmatrix}\right)=-1\\
c_{31}&=\det\left(\begin{smallmatrix}1&0&0\\1&1&0\\0&0&1\end{smallmatrix}\right)=1,\ &
c_{32}&=-\det\left(\begin{smallmatrix}1&0&0\\0&1&0\\-1&0&1\end{smallmatrix}\right)=-1,\ &
c_{33}&=\det\left(\begin{smallmatrix}1&1&0\\0&1&0\\-1&0&1\end{smallmatrix}\right)=1,\ &
c_{34}&=-\det\left(\begin{smallmatrix}1&1&0\\0&1&1\\-1&0&0\end{smallmatrix}\right)=1\\
c_{41}&=-\det\left(\begin{smallmatrix}1&0&0\\1&1&0\\0&1&1\end{smallmatrix}\right)=-1,\ &
c_{42}&=\det\left(\begin{smallmatrix}1&0&0\\0&1&0\\0&1&1\end{smallmatrix}\right)=1,\ &
c_{43}&=-\det\left(\begin{smallmatrix}1&1&0\\0&1&0\\0&0&1\end{smallmatrix}\right)=-1,\ &
c_{44}&=\det\left(\begin{smallmatrix}1&1&0\\0&1&1\\0&0&1\end{smallmatrix}\right)=1
\end{alignat*}
\end{board}
\begin{board}[\(\det A\) along row \(1\) (Theorem~\ref{b5.21}(a)); \(A^{-1}\) (Theorem~\ref{b5.24}):]
\[
\det A=1\cdot c_{11}+1\cdot c_{12}+0\cdot c_{13}+0\cdot c_{14}=2,\qquad
A^{-1}=\frac12\adj A=\frac12\begin{pmatrix}1&-1&1&-1\\1&1&-1&1\\-1&1&1&-1\\1&-1&1&1\end{pmatrix}
\]
\end{board}

\Needspace{13\baselineskip}
\begin{theorem}\label{b5.27}
Let \(A\in\R^{n\times n}\) (\(n\geqslant2\)) with \(\det A\ne0\), \(b\in\R^n\), and \(A_j\) = \(A\) with column \(j\)
replaced by \(b\). Then \(Ax=b\) has exactly one solution \(x\), and \(x_j=\det A_j/\det A\) (\(1\leqslant j\leqslant n\))
(\textbf{Cramer's rule}).
\end{theorem}
\begin{proof}
\begin{board}[\(A^{-1}=\frac1{\det A}\adj A\) (Theorem~\ref{b5.24}); \(A_j(i|j)=A(i|j)\), entry \((i,j)\) of \(A_j\) is \(b_i\):]
\boardtop
\begin{align*}
Ax=b&\implies x=A^{-1}(Ax)=A^{-1}b;\quad A(A^{-1}b)=b&&\text{(\lec{1}, Theorem~3(a))}\\
&\implies\text{the only solution is }x=\tfrac1{\det A}(\adj A)b\\
\bigl((\adj A)b\bigr)_j&=\textstyle\sum_i(\adj A)_{ji}b_i=\sum_ic_{ij}b_i\\
\det A_j&=\textstyle\sum_ib_i(-1)^{i+j}\det A(i|j)=\sum_ic_{ij}b_i&&\text{(Theorem~\ref{b5.15}, column \(j\))}\\
x_j&=\det A_j/\det A
\end{align*}
\end{board}
\nopagebreak
\end{proof}

\Needspace{16\baselineskip}
\begin{example}\label{b5.28}
For \(A\) of Example~\ref{b5.26} and \(b=(3,1,0,0)'\), find \(x_3\) in \(Ax=b\) by Cramer's rule, and check it against
\(x=A^{-1}b\).
\end{example}
\textit{Solution.}
\begin{board}[\(\det A=2\); row \(3\) of \(A_3\) has one nonzero entry, at \((3,4)\); expand along it, then along row
\(3\) (Theorem~\ref{b5.21}(a)); Theorem~\ref{b5.27}; \((\adj A)b\) by \lec{1}, Theorem~4(b):]
\[
\begin{gathered}
\det A_3=\det\begin{pmatrix}1&1&3&0\\0&1&1&0\\0&0&0&1\\-1&0&0&1\end{pmatrix}
=(-1)^{3+4}\det\begin{pmatrix}1&1&3\\0&1&1\\-1&0&0\end{pmatrix}
=-(-1)(-1)^{3+1}\det\begin{pmatrix}1&3\\1&1\end{pmatrix}=-2\\
x_3=\frac{\det A_3}{\det A}=\frac{-2}2=-1\\
x=\tfrac12(\adj A)b=\tfrac12\bigl(3\col_1(\adj A)+\col_2(\adj A)\bigr)=\tfrac12\bigl(3(1,1,-1,1)'+(-1,1,1,-1)'\bigr)=(1,2,-1,1)'
\end{gathered}
\]
\end{board}

\subsection*{5. Determinants of block matrices}

\Needspace{11\baselineskip}
\begin{theorem}\label{b5.29}
Let \(A\in\R^{p\times p}\), \(B\in\R^{p\times q}\), \(C\in\R^{q\times p}\), \(D\in\R^{q\times q}\). Then
\(\det\left(\begin{smallmatrix}A&B\\0&D\end{smallmatrix}\right)=\det A\det D
=\det\left(\begin{smallmatrix}A&0\\C&D\end{smallmatrix}\right)\); in particular \(\det\diag(A,D)=\det A\det D\).
\end{theorem}
\begin{proof}
\begin{board}[Three factors: expand \(p\) times along row \(1\), \(q\) times along the last row (Theorem~\ref{b5.21}(a));
Theorem~\ref{b5.14}(b):]
\[
\det\begin{pmatrix}I_p&0\\0&D\end{pmatrix}=\det D,\qquad\det\begin{pmatrix}A&0\\0&I_q\end{pmatrix}=\det A,\qquad
\det\begin{pmatrix}I_p&B\\0&I_q\end{pmatrix}=1
\]
\end{board}
\begin{board}[The product (Theorem~\ref{b5.3}(c)); Theorem~\ref{b5.18}(b); the transpose (Theorem~\ref{b5.3}(b),
Theorem~\ref{b5.21}(b)):]
\boardtop
\begin{align*}
\begin{pmatrix}I_p&0\\0&D\end{pmatrix}\begin{pmatrix}I_p&B\\0&I_q\end{pmatrix}\begin{pmatrix}A&0\\0&I_q\end{pmatrix}
&=\begin{pmatrix}I_p&0\\0&D\end{pmatrix}\begin{pmatrix}A&B\\0&I_q\end{pmatrix}=\begin{pmatrix}A&B\\0&D\end{pmatrix}
\implies\det\begin{pmatrix}A&B\\0&D\end{pmatrix}=\det D\cdot1\cdot\det A\\
\det\begin{pmatrix}A&0\\C&D\end{pmatrix}&=\det\begin{pmatrix}A'&C'\\0&D'\end{pmatrix}=\det A'\det D'=\det A\det D
\end{align*}
\end{board}
\nopagebreak
\end{proof}

\Needspace{13\baselineskip}
\begin{theorem}\label{b5.30}
Let \(M=\begin{pmatrix}A&B\\C&D\end{pmatrix}\) with \(A\in\R^{p\times p}\), \(B\in\R^{p\times q}\), \(C\in\R^{q\times p}\),
\(D\in\R^{q\times q}\).
\begin{enumerate}
\item If \(A\) is invertible, then \(\det M=\det A\det(D-CA^{-1}B)\).
\item If \(D\) is invertible, then \(\det M=\det D\det(A-BD^{-1}C)\).
\end{enumerate}
\end{theorem}
\begin{proof}
\begin{board}[(a) The left factor of Theorem~\ref{b5.9} (\(n=q\)) is lower triangular with diagonal entries \(1\) (the lines after
Theorem~\ref{b5.21}):]
\[
1\cdot\det M=\det\begin{pmatrix}A&B\\0&D-CA^{-1}B\end{pmatrix}=\det A\det(D-CA^{-1}B)\qquad
\text{(Theorems~\ref{b5.18}(b), \ref{b5.29})}
\]
\end{board}
\begin{board}[(b) The left factor is upper triangular with diagonal entries \(1\) (Theorems~\ref{b5.3}(c),
\ref{b5.14}(b)):]
\[
\begin{pmatrix}I_p&-BD^{-1}\\0&I_q\end{pmatrix}M=\begin{pmatrix}A-BD^{-1}C&B-BD^{-1}D\\C&D\end{pmatrix}
=\begin{pmatrix}A-BD^{-1}C&0\\C&D\end{pmatrix}
\]
\(1\cdot\det M=\det(A-BD^{-1}C)\det D\) (Theorems~\ref{b5.18}(b), \ref{b5.29}).
\end{board}
\nopagebreak
\end{proof}
If \(A\) is invertible, \(\det A\ne0\), so \(M\) is invertible \(\iff\) \(D-CA^{-1}B\) is (by (a) and
Theorem~\ref{b5.18}(a)).

\Needspace{10\baselineskip}
\begin{theorem}\label{b5.32}
Let \(A\in\R^{n\times n}\) and \(T\in\R^{m\times m}\) be invertible, \(U\in\R^{n\times m}\) and \(V\in\R^{m\times n}\).
\begin{enumerate}
\item \(\det(A+UTV)=\det A\det T\det(T^{-1}+VA^{-1}U)\). Hence \(A+UTV\) is invertible \(\iff\) \(T^{-1}+VA^{-1}U\) is
invertible, and then (the \textbf{Sherman--Morrison--Woodbury formula})
\[
(A+UTV)^{-1}=A^{-1}-A^{-1}U(T^{-1}+VA^{-1}U)^{-1}VA^{-1}
\]
\item For \(u,v\in\R^n\), \(A+uv'\) is invertible \(\iff\) \(1+v'A^{-1}u\ne0\), and then\\
\mbox{\((A+uv')^{-1}=A^{-1}-\frac1{1+v'A^{-1}u}A^{-1}uv'A^{-1}\)}.
\end{enumerate}
\end{theorem}
\begin{proof}
(a) Let \(M=\left(\begin{smallmatrix}A&-U\\V&T^{-1}\end{smallmatrix}\right)\). As \((T^{-1})^{-1}=T\) (\lec{2},
Theorem~11(a)), the Schur complements of \(A\) and of \(T^{-1}\) in \(M\) are
\(E=T^{-1}-VA^{-1}(-U)=T^{-1}+VA^{-1}U\) and \(F=A-(-U)TV=A+UTV\).
\begin{board}[Determinants:]
\boardtop
\begin{align*}
\det M&=\det A\det E=\det T^{-1}\det F&&\text{(Theorem~\ref{b5.30}(a),(b))}\\
\det F&=\det A\det T\det E,\quad\det A,\det T\ne0&&\text{(Theorem~\ref{b5.18}(c),(a))}\\
F\text{ invertible}&\iff\det F\ne0\iff\det E\ne0\iff E\text{ invertible}&&\text{(Theorem~\ref{b5.18}(a))}
\end{align*}
\end{board}
\begin{board}[\(E\), \(F\) invertible: the upper left block of \(M^{-1}\), with the blocks \(A\), \(-U\), \(V\), \(T^{-1}\)
of \(M\):]
\boardtop
\begin{align*}
(M^{-1})_{11}&=A^{-1}+A^{-1}(-U)E^{-1}VA^{-1}&&\text{(Corollary~\ref{b5.10}, Schur complement \(E\) of \(A\))}\\
(M^{-1})_{11}&=F^{-1}&&\text{(Theorem~\ref{b5.12b}, Schur complement \(F\) of \(T^{-1}\))}\\
F^{-1}&=A^{-1}-A^{-1}UE^{-1}VA^{-1}&&\text{(\(M^{-1}\) is unique: \lec{2}, Definition~10)}
\end{align*}
\end{board}
\begin{board}[(b) In (a), \(m=1\), \(U=u\), \(T=(1)\), \(V=v'\):]
\(A+UTV=A+uv'\) and \(T^{-1}+VA^{-1}U=(1+v'A^{-1}u)\), a \(1\times1\) matrix: invertible \(\iff\)
\(1+v'A^{-1}u\ne0\) (Theorem~\ref{b5.18}(a), \(\det(a)=a\)), with inverse \(\bigl(1/(1+v'A^{-1}u)\bigr)\) (\lec{2}, Definition~10). As a
\(1\times1\) factor acts as a number (\lec{1}, Theorem~3(b)),
\[
A^{-1}u\bigl(1/(1+v'A^{-1}u)\bigr)v'A^{-1}=\frac1{1+v'A^{-1}u}A^{-1}uv'A^{-1}
\]
\end{board}
\nopagebreak
\end{proof}

\bigskip
\textbf{Practice.} The Exercises and Problems of the textbook, Chapter~5: block products and transposes, block
elimination and the Schur complement, determinants by row reduction and by cofactor expansion, the product rule, the
adjugate and Cramer's rule, determinants of block matrices, and the Sherman--Morrison--Woodbury formula.

\end{document}
